\chapter{Architecture Design}
\label{chap:architecture}

The gaps identified in Chapter~\ref{chap:requirements} shaped the platform architecture
presented here. The design builds on the upstream Time.IO system
(Section~\ref{sec:bg-timeio}) as the infrastructure layer and adds a new application
layer on top. The components fall into three categories by provenance:

\begin{description}
  \item[Inherited (upstream).] Components taken from the UFZ Time.IO platform and
    used without modification: PostgreSQL/TimescaleDB, Mosquitto, MinIO, FROST-Server,
    Keycloak, Grafana, and Flyway.
  \item[Inherited and modified.] Components forked from upstream and extended or fixed
    as part of this thesis: the ingestion workers, the per-thing schema provisioning
    logic, the FROST-Server SQL views (replaced), and the configuration database schema
    (new migrations).
  \item[New.] Components designed and implemented entirely in this thesis: the
    Water Data Platform application layer (FastAPI backend, Next.js frontend, Celery
    worker, GeoServer integration), the alert system, the batch import pipeline,
    and the project-oriented data model.
\end{description}

Section~\ref{sec:arch-overview} lays out the
overall system design and the rationale behind the two-layer split.
Section~\ref{sec:arch-components} defines each component's responsibility boundaries.
The database design, including the per-thing schema pattern and the local STA views,
follows in Section~\ref{sec:arch-database}. Section~\ref{sec:arch-mqtt} details the
MQTT-driven event architecture.

%% ----------------------------------------------------------------
\section{Overall System Architecture}
\label{sec:arch-overview}
%% ----------------------------------------------------------------

The platform is composed of two runtime layers, illustrated in
Figure~\ref{fig:arch-overview}.

\begin{landscape}
\begin{figure}[p]
  \centering
  \begin{tikzpicture}[
    box/.style={draw, thick, minimum width=3.2cm, minimum height=1cm, align=center, font=\small},
    newbox/.style={box, fill=blue!8},
    modbox/.style={box, fill=orange!12},
    layer/.style={draw, thick, rounded corners=3pt, inner sep=14pt},
    arr/.style={-{Stealth[length=5pt]}, thick},
    lbl/.style={font=\scriptsize, fill=white, inner sep=2pt},
    every node/.style={font=\small}
  ]

    %% ═══════════════════════════════════════════════
    %% LAYOUT: 3 rows, 2cm gap between packages
    %%   wdp:  x = -0.5 … 6.5   (center 3,   w=7cm)
    %%   tsm:  x =  8.5 … 15.5  (center 12,  w=7cm)
    %%   gap:  x =  6.5 … 8.5   (clear channel)
    %%   data: x =  0.5 … 14.5  (center 7.5, w=14cm)
    %% ═══════════════════════════════════════════════

    %% ── Researcher ──
    \node[font=\normalsize\bfseries] (user) at (7.5, 1) {Researcher};

    %% ── water-dp (label above-left, away from arrow landing) ──
    \node[layer, fill=blue!4] (wdp) at (3, -3.5) [minimum width=7cm, minimum height=5cm] {};
    \node[font=\small\bfseries, anchor=south west] at (wdp.north west) {Water Data Platform (application layer)};
    \node[newbox] (frontend) at (1.2, -2.2) {Next.js\\Frontend};
    \node[newbox] (api)      at (4.8, -2.2) {FastAPI\\Backend};
    \node[newbox] (celery)   at (1.2, -3.8) {Celery\\Worker};
    \node[newbox] (geo)      at (4.8, -3.8) {GeoServer};

    %% ── tsm-orchestration (label above-left, away from arrow landing) ──
    \node[layer, fill=orange!4] (tsm) at (12, -3.5) [minimum width=7cm, minimum height=5cm] {};
    \node[font=\small\bfseries, anchor=south west] at (tsm.north west)
      {tsm-orchestration (infrastructure layer)};
    \node[box] (frost)   at (10.2, -2.2) {FROST-Server};
    \node[box] (kc)      at (13.8, -2.2) {Keycloak};
    \node[box] (grafana) at (10.2, -3.8) {Grafana};
    \node[modbox] (workers) at (13.8, -3.8) {Ingestion\\Workers};
    \node[box] (cron)    at (12,   -5.2) {Cron Scheduler};

    %% ── Shared Data Layer ──
    \node[layer, label={[font=\small\bfseries]north:Shared Data Layer}]
      (data) at (7.5, -9) [minimum width=14cm, minimum height=2.2cm] {};
    \node[box] (db)    at (4,   -9) {PostgreSQL +\\TimescaleDB};
    \node[box] (minio) at (7.5, -9) {MinIO (S3)};
    \node[box] (mqtt)  at (11,  -9) {Mosquitto\\MQTT Broker};

    %% ── MQTT Device ──
    \node[font=\normalsize\bfseries, align=center] (device) at (15.5, -9) {MQTT\\Device};

    %% ═══════ ARROWS ═══════

    %% 1. Researcher → water-dp (arrow lands on right side of package top)
    \draw[arr] (user.south) -- ++(0,-0.3) -| ($(wdp.north east)!0.3!(wdp.north west)$);
    \node[lbl] at (5.2, 0.5) {HTTPS};

    %% 2. Researcher → tsm (arrow lands on right side of package top)
    \draw[arr] (user.south) -- ++(0,-0.3) -| ($(tsm.north east)!0.3!(tsm.north west)$);
    \node[lbl] at (10.5, 0.5) {HTTPS};

    %% 3. Cross-stack: wdp → tsm (horizontal through gap)
    \draw[arr] (wdp.east) -- node[lbl, above, align=center] {STA REST\\OIDC} (tsm.west);

    %% 4. water-dp → data (straight down)
    \draw[arr] (wdp.south) -- node[lbl, right] {SQL, S3, MQTT} (wdp.south |- data.north);

    %% 5. tsm → data (straight down)
    \draw[arr] (tsm.south) -- node[lbl, left] {SQL, S3, MQTT} (tsm.south |- data.north);

    %% 6. MQTT Device → Mosquitto (horizontal)
    \draw[arr] (device.west) -- node[lbl, above] {MQTT/TLS} (mqtt.east);

    %% ═══════ LEGEND ═══════
    \node[newbox, minimum width=1.2cm, minimum height=0.5cm] (leg1) at (1, -11) {};
    \node[right=2pt of leg1, font=\scriptsize] {New (this thesis)};
    \node[modbox, minimum width=1.2cm, minimum height=0.5cm] (leg2) at (5.5, -11) {};
    \node[right=2pt of leg2, font=\scriptsize] {Inherited \& modified};
    \node[box, minimum width=1.2cm, minimum height=0.5cm] (leg3) at (10.5, -11) {};
    \node[right=2pt of leg3, font=\scriptsize] {Inherited (upstream)};

  \end{tikzpicture}
  \caption{High-level architecture of the Hydro Platform.}
  \label{fig:arch-overview}
\end{figure}
\end{landscape}

\textbf{tsm-orchestration} forms the infrastructure layer (inherited and modified): a Mosquitto MQTT broker,
PostgreSQL with TimescaleDB for time-series storage, FROST-Server for OGC SensorThings~API
access, MinIO for S3-compatible object storage, Keycloak for identity management, Grafana for
visualization, and a set of Python-based ingestion workers. It is a heavily modified
fork of the upstream UFZ Time.IO platform described in Section~\ref{sec:bg-timeio}.
The infrastructure services themselves (database, broker, object store) are used as-is;
the modifications concentrate in the ingestion workers, the SQL views for FROST-Server,
and the database schema migrations.

\textbf{The Water Data Platform} (repository \texttt{water-dp}) sits on top as the application layer (new): a FastAPI backend, a Next.js frontend portal,
a Celery\footnote{Celery distributed task queue: \url{https://docs.celeryq.dev/}} worker for background tasks, GeoServer for geospatial layer management, and Redis\footnote{Redis in-memory data store: \url{https://redis.io/}} as
the Celery message broker. It owns all application-level concepts (projects,
alerts, geospatial layers, sensor activity tracking) absent from the upstream TSM platform.
This entire layer was designed and implemented as part of this thesis.

Keeping the inherited infrastructure separate from the new application logic was
a deliberate choice. The boundary means that upstream TSM updates
can be merged into the fork without conflicting with application-layer code, and that
the data platform can be developed and tested against a running TSM instance without modifying
it.

At runtime, all services communicate over a single Docker bridge network named
\texttt{hydro-platform-net}. Each stack defines its own logical network name
(\texttt{default}, \texttt{water\_shared\_net}, \texttt{tsm\_network}), and the
deployment layer Compose file maps all three names to the same physical network.
This allows services to address each other by container name regardless of which stack
defined them. The Nginx reverse proxy in \texttt{tsm-orchestration} serves as the single
entry point, routing requests to the frontend (\texttt{/}), the API (\texttt{/water-api}),
FROST-Server (\texttt{/FROST-Server}), Grafana (\texttt{/grafana}), GeoServer
(\texttt{/geoserver}), and Keycloak (\texttt{/keycloak}).

%% ----------------------------------------------------------------
\section{Component Responsibilities and Boundaries}
\label{sec:arch-components}
%% ----------------------------------------------------------------

Each component owns a well-defined set of responsibilities. The boundaries are enforced by
the repository structure and the communication protocols: the data platform never writes
directly to TSM-owned database tables, and TSM workers are unaware of application-level
concepts such as projects or alerts.

\subsection{\texttt{tsm-orchestration}}
\label{sec:arch-components-tsm}

The TSM stack owns the following responsibilities:

\begin{itemize}
  \item \textbf{Raw time-series storage.} Per-thing PostgreSQL schemas with TimescaleDB
        hypertables for observation data.
  \item \textbf{Infrastructure provisioning.} The thing-setup worker creates schemas,
        MinIO buckets, MQTT credentials, Grafana dashboards, and FROST-Server context
        files when a new Thing is registered.
  \item \textbf{Data ingestion.} Workers for MQTT real-time ingestion, file-based
        ingestion (CSV, JSON), external API synchronization, and SFTP polling.
  \item \textbf{Quality control.} The run-qaqc worker applies automated QC tests and
        stores quality flags alongside observations.
  \item \textbf{Standards-compliant access.} FROST-Server exposes all observation data
        through the OGC SensorThings~API using the local SQL views described in
        Section~\ref{sec:arch-db-staviews}.
  \item \textbf{Authentication infrastructure.} Keycloak provides OIDC-based identity
        management for both stacks.
\end{itemize}

\subsection{Water Data Platform}
\label{sec:arch-components-waterdp}

The data platform owns all application-level features that were absent from the upstream
TSM platform:

\begin{itemize}
  \item \textbf{Project-oriented data model.} Projects group Things, members, and
        geospatial layers. Members are assigned per-project roles (owner, editor, viewer).
  \item \textbf{Sensor management portal.} The Next.js frontend provides CRUD operations
        for Things, Datastreams, parser configurations, and MQTT device types through
        the FastAPI backend, which proxies the operations to the TSM configuration
        database and FROST-Server.
  \item \textbf{Alert evaluation.} Alert definitions specify threshold, no-data, and
        QA/QC-based rules. The Celery worker evaluates QA/QC alerts in response to MQTT
        events and checks sensor inactivity on a periodic schedule.
  \item \textbf{Geospatial layer management.} GeoServer publishes PostGIS-backed layers
        as WMS and WFS services. The API manages workspaces, datastores, and layer
        publishing through the GeoServer REST~API.
  \item \textbf{Batch import.} CSV files with historical observation data are uploaded
        through the API and processed asynchronously by the Celery worker.
  \item \textbf{Sensor simulation.} A simulated sensor feature generates synthetic
        observation data at configurable intervals, enabling testing without physical
        hardware.
\end{itemize}

%% hydro-deploy is a deployment convenience, not an architectural component.
%% Its description is in Section~\\ref{sec:impl-hydro-deploy} (Implementation chapter).

%% ----------------------------------------------------------------
\section{Database Design}
\label{sec:arch-database}
%% ----------------------------------------------------------------

The platform supports two database deployment variants. In the integrated deployment,
both stacks share a single PostgreSQL instance with schema-level
isolation. In standalone mode, each stack runs its own database server. This section describes
the shared variant, which is the primary deployment target.

\subsection{Schema Organization}
\label{sec:arch-db-shared}

The shared PostgreSQL instance contains three categories of schemas:

\begin{description}
  \item[config\_db schema] Owned by TSM. Stores the configuration database tables that
    drive the ingestion pipeline: \texttt{thing} (sensor definitions), \texttt{project}
    (logical groupings), \texttt{database} (per-project connection credentials),
    \texttt{s3\_store} (MinIO bucket configuration), \texttt{file\_parser} and
    \texttt{file\_parser\_type} (parser definitions), \texttt{mqtt} and
    \texttt{mqtt\_device\_type} (MQTT device configurations), and \texttt{ext\_api} and
    \texttt{ext\_sftp} (external data source definitions). The mapping between a Thing's
    UUID and its schema name is stored in \texttt{public.schema\_thing\_mapping}.

  \item[Per-thing schemas] Owned by TSM. Each onboarded Thing receives a dedicated schema
    (e.g.\ \texttt{user\_abc123}) containing its own \texttt{thing}, \texttt{datastream},
    \texttt{observation}, and \texttt{location} tables. The provisioning process is
    described in Section~\ref{sec:arch-db-perthings}.

  \item[water\_dp schema] Owned by the data platform. Stores application-state tables
    managed by Alembic migrations. The key tables are \texttt{projects} (project
    definitions with Keycloak group linkage), \texttt{project\_sensors} (Thing-to-project
    associations), \texttt{project\_members} (per-project RBAC with owner/editor/viewer
    roles), \texttt{alert\_definitions} and \texttt{alerts} (the alert evaluation
    subsystem), \texttt{geo\_layers} and \texttt{geo\_features} (geospatial data with
    PostGIS geometry columns), \texttt{simulations} (sensor simulation configuration),
    and \texttt{sensor\_activity\_configs} (inactivity monitoring).
\end{description}

The data platform API connects to the shared database with
\texttt{search\_path=water\_dp,public}, which allows it to read TSM configuration tables
in the \texttt{config\_db} schema through explicit schema-qualified queries while defaulting
to the \texttt{water\_dp} schema for its own tables.

\subsection{Standalone Deployment Variant}
\label{sec:arch-db-standalone}

When the data platform is deployed without \texttt{tsm-orchestration}, it starts its own
PostgreSQL instance (\texttt{postgres-app}) with PostGIS. In this mode, the FROST~API client
communicates with an external TSM deployment over HTTP rather than accessing per-thing schemas
directly. The TSM integration overlay (\texttt{docker-compose.tsm.yml}) disables the standalone
database by setting its profile to \texttt{donotstart} and redirects the API's
\texttt{DATABASE\_URL} to the TSM PostgreSQL instance.

\subsection{Per-Thing Schema Provisioning}
\label{sec:arch-db-perthings}

When a new Thing is registered, the thing-setup worker creates a dedicated PostgreSQL schema.
The provisioning sequence is:

\begin{enumerate}
  \item Create a PostgreSQL role with login credentials and a schema owned by that role.
  \item Create the \texttt{observation} table and convert it to a TimescaleDB hypertable
        with time-based chunking. TimescaleDB automatically partitions observations into
        time-range chunks, enabling efficient range queries and background compression
        of older data.
  \item Create the \texttt{datastream}, \texttt{thing}, and \texttt{location} tables.
  \item Deploy the local STA views (Section~\ref{sec:arch-db-staviews}) into the schema.
  \item Grant the FROST-Server database role read access to all views and tables.
\end{enumerate}

This per-thing isolation guarantees that one sensor's data volume or query load does not
affect another, and it simplifies access-control policies at the database level. The mapping
between a Thing's UUID and its schema name is recorded in
\texttt{public.schema\_thing\_mapping}.

\subsection{Local STA Views}
\label{sec:arch-db-staviews}

FROST-Server reads its entity data from PostgreSQL views rather than from a fixed table layout.
The upstream TSM platform defines these views against tables managed by the external SMS
service, which was non-functional (Section~\ref{sec:bg-timeio-limitations}). As part of this
thesis, an alternative set of nine local views was introduced in
\texttt{src/sql/sta\_views\_local/}, one per STA entity type: \texttt{THINGS},
\texttt{DATASTREAMS}, \texttt{OBSERVATIONS}, \texttt{LOCATIONS}, \texttt{THINGS\_LOCATIONS},
\texttt{SENSORS}, \texttt{OBSERVED\_PROPERTIES}, \texttt{FEATURES\_OF\_INTEREST}, and
\texttt{HISTORICAL\_LOCATIONS}.

The views for \texttt{THINGS}, \texttt{DATASTREAMS}, \texttt{OBSERVATIONS}, and
\texttt{LOCATIONS} derive their data from the per-thing schema tables and the configuration
database. For example, the \texttt{DATASTREAMS} view extracts unit-of-measurement metadata
from the JSONB \texttt{properties} column and constructs the STA-compliant JSON structure
that FROST expects. The remaining four entity types (\texttt{SENSORS},
\texttt{OBSERVED\_PROPERTIES}, \texttt{FEATURES\_OF\_INTEREST},
\texttt{HISTORICAL\_LOCATIONS}) have no local equivalent and are provided as empty stub
views so that FROST-Server starts without errors.

A runtime toggle (\texttt{USE\_LOCAL\_STA\_VIEWS}, defaulting to \texttt{true}) controls
which view set is deployed during Thing provisioning. All code related to this workaround
is marked with the label \texttt{TSM-TEMP} in the codebase.

%% ----------------------------------------------------------------
\section{MQTT-Driven Event Architecture}
\label{sec:arch-mqtt}
%% ----------------------------------------------------------------

The Mosquitto MQTT broker serves as the primary integration backbone between all TSM workers
and the data platform Celery worker. All workers follow a message-driven architecture:
each subscribes to one or more topics, performs its designated processing, and optionally
publishes results to downstream topics. This design decouples the components and enables
independent scaling.

\subsection{Topic Structure and Event Flow}
\label{sec:arch-mqtt-topics}

Table~\ref{tab:mqtt-topics} lists the MQTT topics and the workers that produce and consume
them. The Thing creation lifecycle is illustrated in Figure~\ref{fig:mqtt-thing-lifecycle},
and the data ingestion event flow is shown in Figure~\ref{fig:mqtt-data-ingestion}.

\begin{longtable}{p{0.30\textwidth} p{0.30\textwidth} p{0.30\textwidth}}
  \caption{MQTT topics and their producers and consumers.}
  \label{tab:mqtt-topics} \\
  \textbf{Topic} & \textbf{Producer} & \textbf{Consumer} \\
  \hline
  \endfirsthead
  \textbf{Topic} & \textbf{Producer} & \textbf{Consumer} \\
  \hline
  \endhead
  \texttt{frontend\_thing\_update}
    & Data platform API
    & configdb-updater \\[4pt]
  \texttt{configdb\_update}
    & configdb-updater
    & thing-setup \\[4pt]
  \texttt{mqtt\_ingest/\{user\}/\#}
    & MQTT devices
    & mqtt-ingest \\[4pt]
  \texttt{object\_storage\_notification}
    & MinIO (S3 events)
    & file-ingest \\[4pt]
  \texttt{data\_parsed}
    & file-ingest, mqtt-ingest, sync-extapi
    & run-qaqc, data platform worker \\[4pt]
  \texttt{qaqc\_done}
    & run-qaqc
    & data platform worker \\[4pt]
  \texttt{sync\_ext\_apis}
    & cron scheduler
    & sync-extapi \\[4pt]
  \texttt{sync\_ext\_sftp}
    & cron scheduler
    & sync-extsftp \\[4pt]
  \texttt{userinfo\_keycloak}
    & Keycloak (login event)
    & grafana-user-orgs \\[4pt]
  \texttt{monitor\_mqtt}
    & cron scheduler
    & monitor-mqtt \\
\end{longtable}

\begin{landscape}
\begin{figure}[p]
  \centering
  \includegraphics[width=\linewidth,height=0.85\textheight,keepaspectratio]{figures/mqtt-thing-lifecycle.pdf}
  \caption{Sequence diagram: Thing creation lifecycle.}
  \label{fig:mqtt-thing-lifecycle}
\end{figure}
\end{landscape}

\begin{landscape}
\begin{figure}[p]
  \centering
  \includegraphics[width=\linewidth,height=0.85\textheight,keepaspectratio]{figures/mqtt-data-ingestion.pdf}
  \caption{Sequence diagram: data ingestion event flow.}
  \label{fig:mqtt-data-ingestion}
\end{figure}
\end{landscape}

\subsection{Thing Lifecycle}
\label{sec:arch-mqtt-lifecycle}

When a Thing is created or modified through the data platform portal, the API publishes a
configuration message to the \texttt{frontend\_thing\_update} topic. The configdb-updater
worker receives the message, validates the content version, and upserts the configuration into
the \texttt{config\_db} schema. Upon successful storage, it publishes a notification to
\texttt{configdb\_update}, which triggers the thing-setup worker. The thing-setup worker
executes the five-step provisioning sequence described in Section~\ref{sec:arch-db-perthings},
plus MQTT credential creation, MinIO bucket provisioning, and Grafana dashboard setup. This
chain ensures that after a Thing is created through the portal, its entire infrastructure is
provisioned automatically without manual intervention.

\subsection{Alert Evaluation}
\label{sec:arch-mqtt-alerts}

The data platform Celery worker subscribes to two MQTT topics to integrate with the TSM
event flow:

\begin{itemize}
  \item \textbf{\texttt{qaqc\_done}} --- triggers evaluation of QA/QC alert rules for the
        affected Thing. When the run-qaqc worker finishes processing quality flags, the
        alert evaluator queries the observation table for the flagged-to-total ratio within
        a configurable time window. If the ratio exceeds the threshold defined in the alert
        rule, an alert is created with status \texttt{active}.
  \item \textbf{\texttt{data\_parsed}} --- triggers sensor activity recording. The worker
        updates the \texttt{last\_seen\_at} timestamp for the affected Thing, which is
        used by the inactivity monitoring subsystem.
\end{itemize}

The MQTT subscription runs as a daemon thread inside the Celery worker process, started on the
\texttt{worker\_ready} signal. When a message arrives, the handler dispatches a Celery task
for asynchronous processing, ensuring that the MQTT callback does not block.

Sensor inactivity checks remain on a periodic Celery Beat schedule (every 30~minutes). This
hybrid approach (event-driven for QA/QC alerts, periodic for inactivity) was chosen
because inactivity is defined by the \emph{absence} of events, which cannot be detected by
subscribing to a topic.

The alert state machine defines three states: \texttt{active} (rule violation detected),
\texttt{acknowledged} (a user has seen the alert, stored with the acknowledger's identity
and timestamp), and \texttt{resolved} (the condition no longer holds or was manually
resolved). QA/QC alerts are automatically resolved when the flagged ratio drops below
the threshold. Inactivity alerts are resolved when the sensor sends new data.

%% ----------------------------------------------------------------
\section{Deployment Strategy}
\label{sec:arch-deployment}
%% ----------------------------------------------------------------

The deployment is managed by a third repository, \texttt{hydro-deploy}, which has no
runtime presence and is not part of the platform architecture itself. It assembles the
runtime stacks from multiple Docker Compose files and provides a Makefile-based operator
interface. Both stacks can be deployed and operated independently; the deployment layer
simply lets them run as a single unit. This section describes the Compose file merging
strategy, secret management, and Makefile targets.
Rootless Podman is also supported as an alternative container runtime.

\subsection{Compose File Merging}
\label{sec:arch-deploy-compose}

Docker Compose supports layering multiple \texttt{-f} files, where later files override or
extend services defined in earlier ones. The platform uses four Compose files in Docker mode:

\begin{enumerate}
  \item \texttt{tsm-orchestration/docker-compose.yml} --- defines all TSM core services
        (database, MQTT broker, FROST-Server, MinIO, Keycloak, Grafana, Nginx proxy,
        all ingestion workers, and the cron scheduler).
  \item \texttt{water-dp/docker-compose.yml} --- defines the standalone data platform
        services (PostgreSQL with PostGIS, Redis, FastAPI API, GeoServer, Celery worker,
        Next.js frontend).
  \item \texttt{water-dp/docker-compose.tsm.yml} --- the TSM integration overlay. Disables
        the standalone PostgreSQL by setting its profile to \texttt{donotstart}, redirects
        the API and worker \texttt{DATABASE\_URL} to the TSM database, and adds MinIO
        and FROST connection parameters.
  \item \texttt{hydro-deploy/docker-compose.yml} --- the final layer. Maps all logical
        network names to the physical \texttt{hydro-platform-net} bridge network, applies
        startup ordering constraints, and adjusts environment variables.
\end{enumerate}

An optional fifth file (\texttt{docker-compose.tunnel.yml}) adds a Cloudflare Tunnel
container for exposing the platform to the internet without a public IP address.

\subsection{Secret Management}
\label{sec:arch-deploy-secrets}

A single \texttt{.env} file at the \texttt{hydro-deploy} root contains all credentials for
both stacks. The \texttt{generate-secrets.sh} script generates cryptographically secure values
using Python's \texttt{secrets} module. Passwords that must be identical across stacks (e.g.\
the database password used by both TSM services and the data platform API) are generated
once and written to all corresponding variables. These pairs are annotated with \texttt{@sync}
comments in the script. The script is idempotent: it replaces only placeholder values
(\texttt{changeme} or empty strings) and never overwrites previously generated secrets.

\subsection{Makefile Operator Interface}
\label{sec:arch-deploy-makefile}

The \texttt{hydro-deploy} Makefile exposes the following primary targets:

\begin{description}
  \item[make setup] Clones the \texttt{tsm-orchestration} and \texttt{water-dp} repositories,
    creates the \texttt{.env} file from the template, and establishes symlinks for shared
    configuration.
  \item[make secrets] Runs \texttt{generate-secrets.sh} to populate the \texttt{.env} file
    with cryptographically secure passwords.
  \item[make build] Builds locally compiled Docker images (Keycloak themes, data platform
    API and frontend). When \texttt{PREBUILT=1} is set, this step is skipped and
    \texttt{make up} pulls pre-built images from the GitHub Container Registry (GHCR)
    instead, using an additional \texttt{docker-compose.prebuilt.yml} overlay that
    replaces all \texttt{build:} directives with \texttt{image:} references.
    A specific release can be pinned with \texttt{IMAGE\_TAG}
    (e.g.\ \texttt{make PREBUILT=1 IMAGE\_TAG=v1.0 up}).
  \item[make up] Creates the \texttt{hydro-platform-net} network, then starts all services
    in a staged order: TSM core services first, then the GeoServer database, then
    GeoServer, and finally the data platform API, worker, and frontend.
  \item[make migrate] Applies pending Alembic database migrations for the \texttt{water\_dp}
    schema.
  \item[make seed] Populates the platform with test sensors and projects.
  \item[make check] Health-checks all service endpoints and validates that \texttt{@sync}
    credential pairs are consistent.
\end{description}

\subsection{Podman Rootless Support}
\label{sec:arch-deploy-podman}

The platform also supports rootless Podman for environments where the Docker daemon is
unavailable. A preprocessing script (\texttt{podman\_prep.py}) transforms the standard
Compose files into Podman-compatible variants by applying seven automated transformations:
removing unsupported profile declarations, stripping incompatible \texttt{depends\_on}
conditions, converting YAML booleans and environment variable formats, absolutizing
relative paths, and removing \texttt{tmpfs} UID/GID options rejected by rootless
\texttt{crun}. A hand-maintained overlay handles GeoServer UID mapping and network
pre-creation. The Makefile auto-detects the available container engine.
